\editionPage[\pdfbookmark[0]{Corrigés des brevets blancs}{bm-bilan-A-cor-auto-parent}]{bilan-A-cor-auto}{Brevet A · Corrigé}{Automatismes expliqués}{\hyperref[bilan-A-auto]{Énoncé p.\,\pageref*{bilan-A-auto}}}
\begin{blocEdition}{bilan-A-cor-auto}
\begin{questions}
\item $\res8$ : la multiplication est prioritaire.
\item $4/6+1/6=\res{5/6}$.
\item $\res{4{,}5\times10^5}$.
\item $\res{x=-9\text{ ou }x=9}$ : deux solutions opposées.
\item $\res{x^2+6x+9}$.
\item $\res{x\le-4}$ : diviser par $-2$ change le sens.
\item $0{,}15\times80=\res{12}$.
\item $-4+1=\res{-3}$.
\item 7 parts de 5 € : $\res{10\text{ € et }25\text{ €}}$.
\item La troisième valeur est $\res5$.
\item $0{,}5+0{,}4-0{,}2=\res{0{,}7}$.
\item $\res{\vect{AB}=\vect{DC}}$ : côtés opposés de même sens.\end{questions}
\end{blocEdition}
\editionPage{bilan-A-cor-1}{Brevet A · Corrigé 1}{Un atelier de carnets}{\hyperref[bilan-A-ex-1]{Énoncé p.\,\pageref*{bilan-A-ex-1}}}
\begin{blocEdition}{bilan-A-cor-1}
\editionEtape{1 · Partager sans reste}Le nombre de carnets divise 168 et 126. $168=2^3\times3\times7$ et $126=2\times3^2\times7$. Le plus grand diviseur commun est $2\times3\times7=42$. On prépare $\res{42\text{ carnets}}$, chacun contenant $\res4$ feuilles bleues et $\res3$ blanches.\editionEtape{2 · Recettes et solde}Les recettes sont $42\times4{,}50=189$ €. Le solde, hors autres frais, est $189-78=\res{111\text{ €}}$.\editionEtape{3 · Une comparaison de prix}Diminuer de 10 \% revient à multiplier par 0,9 : $5\times0{,}9=\res{4{,}50\text{ €}}$. Les prix unitaires sont égaux.\editionRetenir{Le maximum de lots identiques correspond au plus grand diviseur commun, pas au plus petit multiple commun.}
\end{blocEdition}
\editionPage{bilan-A-cor-2}{Brevet A · Corrigé 2}{Une voile divisée}{\hyperref[bilan-A-ex-2]{Énoncé p.\,\pageref*{bilan-A-ex-2}}}
\begin{blocEdition}{bilan-A-cor-2}
\editionEtape{1 · Le triangle rectangle}Le théorème de Pythagore donne $BC^2=12^2+9^2=225$, donc $\res{BC=15\text{ m}}$.\editionEtape{2 · La réduction}Les alignements et le parallélisme permettent d’appliquer le théorème de Thalès : $AM/AB=AN/AC=MN/BC=4/12=1/3$. Ainsi $\res{AN=3\text{ m}}$ et $\res{MN=5\text{ m}}$.\editionEtape{3 · Une translation}Par définition de la translation de vecteur $\vect{MB}$, $\vect{ND}=\vect{MB}$ : $\res{MBDN\text{ est un parallélogramme}}$. Donc $(BD)\parallel(MN)$. Comme $(MN)\parallel(BC)$, les droites (BD) et (BC) sont parallèles et ont le point B en commun : elles sont confondues, et $\res{D\in(BC)}$.\editionRetenir{Deux droites parallèles ayant un point commun sont confondues : c’est un argument classique pour prouver un alignement.}\label{fig-bilan-A-cor-2}\figureGeo[corrige]{brevet-A-geo}

\end{blocEdition}
\editionPage{bilan-A-cor-3}{Brevet A · Corrigé 3}{Deux locations de matériel}{\hyperref[bilan-A-ex-3]{Énoncé p.\,\pageref*{bilan-A-ex-3}}}
\begin{blocEdition}{bilan-A-cor-3}
\editionEtape{1 · Quatre journées}A coûte $18+28=\res{46\text{ €}}$ ; B coûte $\res{40\text{ €}}$. B est moins chère de 6 €.\editionEtape{2 · Le seuil}$18+7n=10n$ donne $3n=18$, donc $\res{n=6}$ ; les deux prix valent 60 €. L’inéquation $18+7n<10n$ équivaut à $18<3n$, soit $n>6$ : $\res{A\text{ est moins chère à partir de 7 jours}}$, B l’est pour moins de 6 jours.\editionEtape{3 · Un budget}Chez A, $18+7n\le90$ équivaut à $n\le72/7\approx10{,}3$ : maximum $\res{10\text{ jours}}$. Chez B, $10n\le90$ équivaut à $n\le9$ : maximum $\res{9\text{ jours}}$. On ne garde que des nombres entiers de jours.\label{fig-bilan-A-cor-3}\figureGeo[corrige]{brevet-A-tarifs}

\end{blocEdition}
\editionPage{bilan-A-cor-4}{Brevet A · Corrigé 4}{Des scores et un tirage}{\hyperref[bilan-A-ex-4]{Énoncé p.\,\pageref*{bilan-A-ex-4}}}
\begin{blocEdition}{bilan-A-cor-4}
\editionEtape{1 · La série}La somme pondérée vaut $0+5+16+15=36$, donc la moyenne vaut $\res{1{,}8}$. Les effectifs cumulés sont 2, 7, 15, 20. Les rangs 10 et 11 valent 2 : médiane $\res2$. $Q_1$ est la 5\textsuperscript{e} valeur : $\res{Q_1=1}$ ; $Q_3$ est la 15\textsuperscript{e} : $\res{Q_3=2}$.\editionEtape{2 · Deux rouges}Le tableau compte $4\times4=16$ couples équiprobables, dont $3\times3=9$ formés de deux rouges : $\res{9/16}$.\editionEtape{3 · Une bleue, au moins une rouge}Exactement une bleue : $3+3=6$ couples, soit $\res{6/16=3/8}$. Au moins une rouge est le contraire de « deux bleues », qui ne contient que le couple $(B\,;B)$ : $\res{1-1/16=15/16}$.\label{fig-bilan-A-cor-4}\figureGeo[corrige]{brevet-A-tableau}

\end{blocEdition}
\editionPage{bilan-B-cor-auto}{Brevet B · Corrigé}{Automatismes expliqués}{\hyperref[bilan-B-auto]{Énoncé p.\,\pageref*{bilan-B-auto}}}
\begin{blocEdition}{bilan-B-cor-auto}
\begin{questions}
\item $\res6$.
\item $3/5\times10/9=\res{2/3}$.
\item $\res{10^{-2}}$.
\item $\res{2^2\times3\times5}$.
\item $\res{(x-7)(x+7)}$.
\item $\res{x=3\text{ ou }x=-1/2}$.
\item $50\times1{,}12=\res{56}$.
\item $-3x=12$, donc $\res{-4}$.
\item $\frac43\pi\times27=\res{36\pi\text{ cm}^3}$.
\item $8\times0{,}25=2$ : la 2\textsuperscript{e} valeur, $\res3$.
\item $1-0{,}3=\res{0{,}7}$.
\item $\res{\vect{AC}}$, d’après la relation de Chasles.\end{questions}
\end{blocEdition}
\editionPage{bilan-B-cor-1}{Brevet B · Corrigé 1}{Un abri et sa pente}{\hyperref[bilan-B-ex-1]{Énoncé p.\,\pageref*{bilan-B-ex-1}}}
\begin{blocEdition}{bilan-B-cor-1}
\editionEtape{1 · Longueur du versant}$BC^2=4^2+2{,}5^2=22{,}25$, donc $BC=\sqrt{22{,}25}\approx\res{4{,}72\text{ m}}$.\editionEtape{2 · Angle de pente}$\tan\widehat B=2{,}5/4=0{,}625$. La fonction inverse donne $\widehat B\approx32{,}005^\circ$, soit $\res{32^\circ}$.\editionEtape{3 · Achat entier}L’aire utile vaut $6\sqrt{22{,}25}\approx28{,}3019$ m$^2$. Avec les découpes, on multiplie par 1,08 : environ 30,5660 m$^2$. Le nombre de plaques doit être au moins $30{,}5660/2\approx15{,}283$. Il faut $\res{16\text{ plaques}}$. Quinze ne couvriraient que 30 m$^2$.\label{fig-bilan-B-cor-1}\figureGeo[corrige]{brevet-B-geo}

\end{blocEdition}
\editionPage{bilan-B-cor-2}{Brevet B · Corrigé 2}{Une bouteille et une boule}{\hyperref[bilan-B-ex-2]{Énoncé p.\,\pageref*{bilan-B-ex-2}}}
\begin{blocEdition}{bilan-B-cor-2}
\editionEtape{1 · Capacité}Le volume vaut $\pi\times4^2\times20=\res{320\pi\text{ cm}^3}$. Un cm$^3$ vaut un mL : capacité $\res{1005\text{ mL}}$ au mL.\editionEtape{2 · Remplissage}Le volume visé vaut $0{,}9\times320\pi=288\pi$ mL. La durée vaut $288\pi/50\approx\res{18{,}1\text{ s}}$.\editionEtape{3 · Une boule}Le récipient a un volume de $\frac43\pi\times7^3=\frac{1372\pi}3\approx1436{,}8$ cm$^3$, supérieur à $1005$ cm$^3$ : $\res{\text{il est assez grand}}$. La part remplie vaut $\dfrac{320\pi}{1372\pi/3}=\dfrac{240}{343}\approx\res{70\,\%}$.\label{fig-bilan-B-cor-2}\figureGeo[corrige]{brevet-B-cylindre}

\end{blocEdition}
\editionPage{bilan-B-cor-3}{Brevet B · Corrigé 3}{Un script de figures}{\hyperref[bilan-B-ex-3]{Énoncé p.\,\pageref*{bilan-B-ex-3}}}
\begin{blocEdition}{bilan-B-cor-3}
\editionEtape{1 · Une variable à chaque tour}Les carrés ont pour côtés $\res{10,15,20,25}$ unités : l’ajout de 5 vient après le tracé. Après le quatrième tour, c vaut 30 et la condition $c>25$ devient vraie : $\res{4\text{ passages}}$.\editionEtape{2 · Longueur dessinée}La somme des côtés vaut 70. Les périmètres totalisent $4\times70=\res{280}$ unités.\editionEtape{3 · État final}Le script affiche $\res{30}$. Les déplacements sans tracer valent $c+5$ à chaque tour, soit 15, 20, 25, 30 : leur somme vaut $\res{90}$ unités. La boucle s’arrête dès que $c>25$ ; aucun appel carré(30) n’est exécuté.\label{fig-bilan-B-cor-3}\figureGeo[corrige]{brevet-B-script}

\end{blocEdition}
\editionPage{bilan-B-cor-4}{Brevet B · Corrigé 4}{Un jeu à deux étapes}{\hyperref[bilan-B-ex-4]{Énoncé p.\,\pageref*{bilan-B-ex-4}}}
\begin{blocEdition}{bilan-B-cor-4}
\editionEtape{1 · Dénombrer}Le tableau donne $4\times2=\res{8\text{ issues équiprobables}}$. Les issues gagnantes sont A1-rouge, A2-rouge, A3-rouge et B-bleue : $P(G)=4/8=\res{1/2}$.\editionEtape{2 · Fréquence et probabilité}$\res{\text{Non}}$ : sur une série de 100 parties, le nombre de gains varie d’une série à l’autre. La probabilité $1/2$ indique que la fréquence des gains se stabilise autour de 0,5 quand le nombre de parties devient très grand.\editionEtape{3 · Réunion}$P(E)=6/8=3/4$, $P(F)=4/8=1/2$ et $E\cap F$ contient 3 issues : $P(E\cap F)=3/8$. La relation du cours donne $P(E\cup F)=3/4+1/2-3/8=\res{7/8}$. Contrôle : seule l’issue B-bleue n’est ni dans E ni dans F.\label{fig-bilan-B-cor-4}\figureGeo[corrige]{brevet-B-tableau}

\end{blocEdition}
